Um zu zeigen, dass $V = \{x \in E : h(x) = 0\}$ ein Unterraum von $E$ ist, überprüfen wir die Unterraumeigenschaften:

Zunächst zeigen wir, dass das Null-Element in $V$ liegt. Da $h$ eine Halbnorm ist, gilt $h(0) = 0$. Somit ist $0 \in V$.

Für die Abgeschlossenheit unter Addition nehmen wir $x, y \in V$ an. Dann gilt $h(x) = 0$ und $h(y) = 0$. Da $h$ subadditiv ist, folgt:

\begin{align*}
h(x + y) &\leq h(x) + h(y) = 0 + 0 = 0.
\end{align*}

Da $h(x + y) \geq 0$ für alle $x, y \in E$, muss $h(x + y) = 0$ gelten, also ist $x + y \in V$.

Für die Abgeschlossenheit unter skalaren Multiplikationen sei $x \in V$ und $\lambda \in \mathbb{R}$. Dann gilt $h(x) = 0$, und aufgrund der positiven Homogenität der Halbnorm folgt:

\begin{align*}
h(\lambda x) &= |\lambda| h(x) = |\lambda| \cdot 0 = 0.
\end{align*}

Daraus folgt, dass $\lambda x \in V$. Somit ist $V$ ein Unterraum von $E$.

Nun definieren wir die Norm $\|\cdot\|_h$ auf dem Faktorraum $E / V$ durch $\|x + V\|_h = h(x)$ für $x \in E$. Wir müssen zeigen, dass diese Definition wohldefiniert ist und die Normaxiome erfüllt.

Die Wohldefiniertheit folgt daraus, dass, wenn $x + V = y + V$, dann $x - y \in V$, also $h(x - y) = 0$. Wegen der Subadditivität der Halbnorm gilt:

\begin{align*}
h(x) &= h(y + (x - y)) \leq h(y) + h(x - y) = h(y) + 0 = h(y).
\end{align*}

Analog folgt $h(y) \leq h(x)$, also $h(x) = h(y)$. Daher ist $\|x + V\|_h = h(x)$ wohldefiniert.

Die Normaxiome sind wie folgt erfüllt:

1. **Positive Definitheit**: $\|x + V\|_h = 0$ impliziert $h(x) = 0$, also $x \in V$, was bedeutet, dass $x + V = V$ das Null-Element im Faktorraum ist.

2. **Positive Homogenität**: Für $\lambda \in \mathbb{R}$ gilt:

\begin{align*}
\|\lambda(x + V)\|_h &= \|\lambda x + V\|_h = h(\lambda x) = |\lambda| h(x) = |\lambda| \|x + V\|_h.
\end{align*}

3. **Subadditivität**: Für $x, y \in E$ gilt:

\begin{align*}
\|x + V + y + V\|_h &= \|x + y + V\|_h = h(x + y) \leq h(x) + h(y) = \|x + V\|_h + \|y + V\|_h.
\end{align*}

Damit ist $\|\cdot\|_h$ eine Norm auf dem Faktorraum $E / V$.

%CHECK: Die Lösung wurde überprüft und keine Schwächen oder Verbesserungsvorschläge wurden im Review identifiziert. Die Herleitung wurde beibehalten und ist korrekt.